\uuid{p59L}
\exo7id{3733}
\titre{exo7 3733}
\auteur{quercia}
\organisation{exo7}
\datecreate{2010-03-11}
\isIndication{false}
\isCorrection{false}
\chapitre{Espace euclidien, espace normé}
\sousChapitre{Projection, symétrie}
\cpge{mpsi-mp2i.espaces-prehilbertiens.a, pcsi.espaces-prehilbertiens.a, ptsi.geometrie-plan.b, mpsi-mp2i.espaces-prehilbertiens.b, pcsi.espaces-prehilbertiens.b, ptsi.geometrie-espace.b, mpsi-mp2i.applications-lineaires.a, pcsi.applications-lineaires.a, ptsi.applications-lineaires.a, bcpst1.geometrie.f, bcpst1.applications-lineaires-matrices, mp-mpi.endomorphismes-euclidien.c, pc.endomorphismes-euclidiens.a, psi.endomorphismes-euclidien.a, pt.isometries-euclidien.a, bcpst2.produit-scalaire.a, bcpst2.applications-lineaires-matrices.a}
\module{Algèbre}
\niveau{L2}
\difficulte{}

\contenu{
\texte{
Soit $E$ un espace euclidien de dimension $n$.
Soit $f : E \to E$ une application non nécéssairement linéaire.
}
\begin{enumerate}
    \item \question{On suppose que $f$ conserve le produit scalaire.
    Démontrer que $f$ est linéaire.}
    \item \question{On suppose que $f$ conserve les distances, c'est à dire :
    $\forall\ \vec x,\vec y \in E,\ \| f(\vec x) - f(\vec y) \| = \|\vec x - \vec y\,\|$.
    Démontrer que $f = f(\vec0) + g$, avec $g \in {\cal O}(E)$.}
\end{enumerate}
}
